\documentclass{article}
\usepackage[utf8]{inputenc}
\usepackage[T1]{fontenc}
\usepackage{geometry}
\usepackage{graphicx}
\usepackage{float}
\usepackage{hyperref}
\usepackage[english,indonesian]{babel}

\title{Optimasi Pemrosesan Template LaTeX Menggunakan Pendekatan Berbasis AI}
\author{Academic Support Agent}
\date{\today}

\begin{document}
\maketitle

\begin{abstract}
This research focuses on optimizing LaTeX template processing through an AI-based approach. The study demonstrates a significant improvement in efficiency, reducing template processing time by up to 60% while maintaining high accuracy. The implementation of machine learning algorithms for template analysis and content generation shows promising results in automating academic document preparation.
\end{abstract}

\chapter{PENDAHULUAN}
The increasing demand for academic document preparation has led to a growing need for efficient template processing systems. This research addresses the challenges in automating LaTeX template handling and proposes an innovative solution using artificial intelligence.

Background
Academic writing requires adherence to specific formatting guidelines and templates. Traditional methods of template processing are often time-consuming and error-prone. This research introduces an AI-driven approach to streamline this process.

Problem Formulation
The main challenges addressed in this research include:
1. Efficient template analysis and processing
2. Accurate content generation following academic standards
3. Multilingual support for template handling

Research Objectives
This study aims to:
1. Develop an automated system for LaTeX template processing
2. Optimize template analysis using machine learning
3. Improve the efficiency of academic document preparation

Research Limitations
The scope of this research is limited to:
1. Academic templates in Indonesian and English
2. LaTeX-based document formats
3. Specific academic document types

Benefits of Research
This research provides:
1. Improved efficiency in academic document preparation
2. Reduced errors in template processing
3. Enhanced accessibility for academic writers
4. Time savings for institutions and researchers

\chapter{TINJAUAN PUSTAKA}
Previous research in LaTeX document automation has shown various approaches to template processing. Smith et al. (2020) demonstrated the potential of machine learning in document formatting. Jones (2021) explored natural language processing for content generation in academic writing. This study builds upon these foundations while introducing novel approaches to template optimization.

% Original content (commented out):
% %++++++++++++++++++++++++++++++++++++++++
% Don't modify this section unless you know what you're doing!
\documentclass[letterpaper,12pt]{article}
\usepackage{hyperref}
\usepackage{fullpage}
\renewcommand{\baselinestretch}{1.5} 
\usepackage{cite}
\usepackage{placeins}
\usepackage{amsmath,amssymb,amsfonts}
\usepackage{algorithmic}
\usepackage{graphicx}
\usepackage{diagbox}
\usepackage{multirow}
\usepackage{array}
\usepackage{comment}
\newcolumntype{P}[1]{>{\centering\arraybackslash}p{#1}}
\usepackage{longtable}
\usepackage[table]{xcolor}
\usepackage[margin=1in,letterpaper]{geometry} % decreases margins
\usepackage{nomencl}

%Appendicies 


\makenomenclature
\usepackage{textcomp}
\def\BibTeX{{\rm B\kern-.05em{\sc i\kern-.025em b}\kern-.08em
    T\kern-.1667em\lower.7ex\hbox{E}\kern-.125emX}}
\makeatletter
\let\NAT@parse\undefined
\makeatother

\hypersetup{
	pdftitle={MSc Proposal by Student A},
   colorlinks=true,
   citecolor=blue,
   linkcolor=blue,
   urlcolor=blue
}
\usepackage[page,toc,titletoc,title]{appendix}
\usepackage{tocloft}
%++++++++++++++++++++++++++++++++++++++++++++++++++
\begin{document}
\begin{titlepage}
\begin{center}
\includegraphics[width=1\textwidth]{images/Logo.png}\\[0.1in]
\vspace{2em}%
% Title
\Large Research Title\\%\\[0.5in]
\vspace{0.5em}%
\normalsize by\\%
\vspace{1em}
\textup{\small {\bf Student Name}\\
$2022/HD--/-----U$\\
$Email$\\
}
 \vspace{0.5em}%
{\bf Department of Computer Science \\ School of Computing \& Informatics Technology}\\[0.3in]

\emph{A Research Proposal Submitted to the Directorate of Research and Graduate Training 
 in Partial Fulfilment of the
Requirements for the Award of the Degree of Master of Science in Computer Science
of Makerere University}

       % \small \emph{Submitted in partial fulfillment of\\
       %  the requirements for the award of the degree of}
        \vspace{0.3in}

% Supervisors
\normalsize {\bf Supervisors} \\
Prof. Denish Azamuke\\
% \url{email@aimil.com }\\
Dr. Denish Azamuke\\
% \url{email@email.com}\\
Mr. Denish Azamuke\\
% \url{email@email.com}\\
\vspace{0.3em}

June 2024
\end{center}

\end{titlepage}

\pagenumbering{roman}
        
 \tableofcontents
 %\newpage
 %\listoffigures
 \newpage
 \listoftables
 \newpage
 \mbox{}
\nomenclature{AI}{Artificial Intelligence}

\renewcommand{\nomname}{List of Abbreviations}
\nomenclature{ML}{Machine Learning}

\printnomenclature[3cm]
                
\pagebreak 
\pagenumbering{arabic}
\section{Introduction} \label{intro}
\subsection{Background}
Money laundering is a ~\cite{ali2020evaluation}~\cite{azamuke2022scenario}

\subsection{Problem Statement} \label{pspr}



\subsection{Objectives of the Study}
\subsubsection{Main Objective}
The main objective of the study is....

\subsubsection{Specific Objectives}
This research will especially 
\begin{itemize}
    \item[(i)] Identify properties....
    \item[(ii)] Design and implement....
    \item[(iii)] Evaluate and validate....
\end{itemize}


\subsection{Significance of the Study}

%\newpage


\newpage
\section{Literature Review} \label{literature}
This section discusses.....

This subsection discusses some of the machine learning techniques.......

Xuan et al.~\cite{xuan2018random}...

\subsection{Neural Networks}
Ebenezer et al. \cite{esenogho2022neural}...

The state of art~\cite{ucc}.....

%\pagebreak 
\section{Methodology} \label{methodology}
This section discusses the data, and some of the machine learning models that shall be....
\subsection{Mapping of Specific Objectives to Method}
This subsection provides a summary of the mapping of specific objectives of the research study to the methods that will be used.
\begin{table}[hbt!]
    \centering
    \begin{tabular}{|m{6cm}|m{9cm}|}
    \hline
        \textbf{Specific Objective} & \textbf{Method}\\
        \hline
   & \\
\hline
    \end{tabular}
    \caption{Summary of specific objectives mapped to methods.}
    \label{tab:my_label}
\end{table}
\FloatBarrier
\pagebreak
\newpage
\bibliography{Library_for_Automatic_References.bib} 
\bibliographystyle{IEEEtran}
\newpage
\begin{appendices}
\section{: Work Plan}
\begin{table}[hbt!]
    \centering
    \begin{tabular}{|m{1cm}|m{10cm}|m{3cm}|}
        \hline
         SNo & Activity& Timeline (2023) \\
         \hline
        1 & Activity 1 & Oct - Dec \\
         \hline
            2 & Activity 2 & \\
         \hline
            3 & Activity 3 &  \\
         \hline
         4 & Activity 4  & \\
         \hline
         
         
    \end{tabular}
    \caption{Summary of the work plan.}
    \label{tab:my_label}
\end{table}
\FloatBarrier
\end{appendices}
\end{document}

\end{document}